\section{Discussion}
\label{sec:discussion}

\subsection{The Shallow Persona Hypothesis}
\label{sec:shallow}

Our results support what we call the \emph{shallow persona hypothesis}: for frontier language models, all practically relevant personas are accessible with minimal context. The persona name alone---supplied via a short system prompt---is sufficient to activate any of the 25 tested personas with near-perfect accuracy. Additional in-context demonstrations provide negligible behavioral improvement ($<$0.4\% mean gain) and only marginal generation quality improvement (+0.18/10).

This finding contrasts with the intuition that more context should unlock deeper or more nuanced persona modes. The gap between $K{=}0$ and $K{=}3$ that \citet{choi2024picle} observed on Llama-2-7B has effectively closed for frontier models: \gptmini already internalizes the mapping from persona names to persona-consistent behavior during pretraining and instruction tuning. The ``context requirement'' for persona elicitation appears to decrease with model capability.

\subsection{Implications for the Total Persona Count}
\label{sec:implications}

If every persona can be activated with $K \leq 3$ examples (or even just a name), the number of accessible personas is approximately the number of distinct persona descriptions that fit in a short prompt. This space is combinatorially large---millions of multi-trait persona profiles can be specified by combining personality, ideology, expertise, and demographic descriptors---but it is \emph{shallow}: no hidden depths require extensive conditioning to discover.

This is consistent with the Linear Representation Hypothesis as applied to personas. If personas correspond to linear directions in activation space \citep{chen2025personavectors, wang2025geometry}, and these directions can be specified by short descriptions, then the persona space is a finite-dimensional manifold navigable by concise natural-language coordinates. The orthogonality of Big Five trait subspaces \citep{wang2025geometry} further suggests that the effective persona dimensionality is modest, even if the combinatorial surface is vast.

\subsection{Why Are All Personas Shallow?}
\label{sec:why_shallow}

Several factors contribute to the absence of context-heavy personas:

\para{Model capability.}
\gptmini is a frontier model trained on massive corpora that include extensive coverage of personality frameworks, ethical theories, political positions, and cultural perspectives. It has already internalized strong priors for all tested persona categories, making zero-shot persona adoption trivial.

\para{Task simplicity.}
The Anthropic benchmark uses binary agree/disagree statements that are often transparent about what they test. A statement like ``I enjoy causing pain to others'' is easily classifiable as consistent with a psychopathy persona regardless of context depth. More nuanced persona expressions may require more context.

\para{System prompt effectiveness.}
Even the $K{=}0$ condition includes a system prompt naming the persona. For a model trained with instruction tuning, the persona name alone provides sufficient specification to activate the appropriate behavioral mode.

\para{RLHF as the binding constraint.}
The only personas showing difficulty are AI-risk ones, where safety training creates resistance. This resistance is about overcoming alignment guardrails, not about accessing latent persona representations. The distinction matters: these are not genuinely ``deep'' personas but rather personas that conflict with the model's trained safety objectives.

\subsection{Limitations}
\label{sec:limitations}

\para{Persona granularity.}
The Anthropic benchmark tests coarse-grained personas defined by single traits or viewpoints. More complex composite personas---for instance, ``a Victorian-era scholar who is secretly sympathetic to anarchism''---might require more context to establish. Our results should be interpreted as applying to the granularity level of the benchmark, not to arbitrarily complex persona specifications.

\para{Ceiling effects.}
\gptmini may be too capable for this benchmark, creating ceiling effects that mask genuine context dependencies. Testing on weaker models (\eg open-source models in the 7B--13B range) might reveal context-dependent personas that frontier models have already internalized.

\para{LLM-as-judge limitations.}
The distinctiveness scores in Experiment~2 are generated by the same model family (\gptmini), introducing potential self-evaluation bias. The judge may systematically over- or under-rate certain persona categories, and its ratings may not correlate with human judgments of persona quality.

\para{Limited $K$ range for generation.}
We test up to $K{=}100$ for generation distinctiveness. Testing $K{=}1000$+ might reveal different patterns, though prompt length limits and cost constrain this exploration.

\para{Single model.}
All results are specific to \gptmini. Different architectures, training approaches, or model scales might exhibit different persona-context relationships. The shallow persona hypothesis may not hold for smaller or less capable models.

\para{No mechanistic analysis.}
We measure behavioral outputs only. Mechanistic analysis---tracking persona vector projections in activation space as a function of $K$---might reveal internal representation changes that are not captured by action consistency or distinctiveness scores.
